% ECONOMICS_PROBLEM: {"id": "I13-621", "title": "Universal coverage and financial protection", "jel_code": "I13", "source_id": "OP-0449"}
% !TeX program = xelatex
\documentclass[11pt,a4paper]{article}
\usepackage[margin=23mm]{geometry}
\usepackage{fontspec}
\setmainfont{Latin Modern Roman}
\usepackage{amsmath,amssymb,mathtools}
\usepackage{xcolor,enumitem,booktabs,longtable,array,xurl,hyperref}
\definecolor{muted}{HTML}{53636C}
\hypersetup{colorlinks=true,urlcolor=blue,linkcolor=blue}
\setlength{\parindent}{0pt}
\setlength{\parskip}{5pt}
\newcommand{\R}{\mathbb R}
\newcommand{\E}{\mathbb E}
\newcommand{\N}{\mathbb N}
\newcommand{\ind}{\mathbf 1}
\newcommand{\Var}{\operatorname{Var}}
\newcommand{\Cov}{\operatorname{Cov}}
\newcommand{\supp}{\operatorname{supp}}
\DeclareMathOperator*{\argmin}{arg\,min}
\DeclareMathOperator*{\argmax}{arg\,max}
\newcommand{\entrymeta}[3]{{\sffamily\small\color{muted}\textbf{\texttt{#1}}\quad|\quad #2\par #3\par}}
\newcommand{\Problemref}[1]{\texttt{#1}}
\newcommand{\JELref}[2]{\texttt{#2} (JEL \texttt{#1})}
\begin{document}
\section*{Universal coverage and financial protection}
\textbf{Problem ID:} \texttt{I13-621}\quad \textbf{JEL:} \texttt{I13}\par
% BEGIN ECONOMICS STATEMENT
\label{problem:OP-0449}
{\sffamily\small\color{muted}Original subject group: Health economics\par}
\label{definitions:problem:30007731}
\textbf{Audit classification:} Published research agenda. Checked source identity and appropriate agenda/background support; expanded intervention, units, population, definedness and causal restrictions. Exact targets and variants are editorial.
\subsection*{Definitions and precise research target}
\textit{Editorial formalization.} Consider uninsured low-income households in one country introducing publicly financed essential health coverage. Let \(a=1\) denote eligibility for the specified benefits package and \(a=0\) the previous financing arrangement. For household \(i\) over five years, let \(X_i(a)\) be out-of-pocket medical spending, \(C_i(a)>0\) total consumption, \(D_i(a)\) medical debt, and \(H_i(a)\) health. For threshold \(\kappa\in(0,1)\), define
\[\tau_F=E[\mathbf1\{X_i(1)/C_i(1)>\kappa\}-\mathbf1\{X_i(0)/C_i(0)>\kappa\}],\qquad\tau_D=E[D_i(1)-D_i(0)],\]
where \(\mathbf1\) is the indicator function. Identify these contrasts using randomized eligibility or a justified phased-rollout design, and report the companion health contrast \(E[H_i(1)-H_i(0)]\). Include people who forgo care, since low spending can reflect unmet needs rather than protection. Account for provider availability, informal charges, taxes financing coverage, and household borrowing responses. Clarify whether estimates apply to eligibility, enrollment, or actual treatment receipt; the proposed target is eligibility.

Use a common five-year window and define \(X_i(a)\) and \(C_i(a)\) as nonnegative real expenditure totals discounted with the same convention, \(C_i(a)>0\) almost surely. State whether total consumption includes medical consumption; the ratio threshold is otherwise ambiguous. \(D_i(a)\) is outstanding medical debt at the end of year five, including a stated interest treatment. Declare the exact benefit package, copayments, provider availability, and financing rule; eligibility does not guarantee enrollment. The indicator contrast lies in \([-1,1]\), while debt effects require finite means. Consumption changes in the denominator are part of the eligibility effect. Report unmet-care and health outcomes jointly. If premium/transport costs are omitted from \(X\), they are separate financial-protection outcomes rather than silently included in out-of-pocket spending. All expectations use a stated baseline population and probability law. Identification means every admissible data-generating model with the same observed-data law yields the same displayed target; otherwise report the identified set. Exact equations, horizons and assumptions are editorial, rather than source-given canonical conjectures.
\subsection*{Variants and assumption changes}
\begin{enumerate}
\item \textbf{Editorial financial-loss variant.} Estimate \(E[X_i(a)+D_i(a)+\operatorname{premiums}_i(a)+\operatorname{travel}_i(a)]\), excluding debt repayments already counted in spending. It differs from catastrophic-spending prevalence.
\item \textbf{Editorial receipt variant.} Distinguish eligibility, enrollment and service use. Identify offer effects directly and receipt effects only with declared exclusion/instrument conditions.
\end{enumerate}
\subsection*{References and source audit}
\textbf{1.} Explicit financial-protection effectiveness and measurement agenda; five-year eligibility estimand is editorial. Locator: BMJ Open 12(3), e052041; Abstract and Results, evidence-gap discussion. Primary indexed publisher passages checked; direct full text failed and PMC returned a challenge page. URL: \url{https://doi.org/10.1136/bmjopen-2021-052041}.
\begin{enumerate}\raggedright
\item\hypertarget{definitions:source:30007731:1}{} Dominika Bhatia; Sujata Mishra; Abirami Kirubarajan; Bernice Yanful; Sara Allin; Erica Di Ruggiero. 2022. \emph{Identifying priorities for research on financial risk protection to achieve universal health coverage: a scoping overview of reviews}. BMJ Open 12(3), e052041; Abstract and Results, evidence-gap discussion.
\url{https://doi.org/10.1136/bmjopen-2021-052041}.
DOI: \url{https://doi.org/10.1136/bmjopen-2021-052041}.
\end{enumerate}

\subsection*{Common conventions: Source mathematical conventions}

These conventions apply to each entry unless explicitly replaced.

\textbf{Models and identification.} A primitive model \(m\) specifies all probability laws, feasible choices, information, equilibrium and measurement rules used by its target. The admissible class \(\mathcal M\) is fixed independently of outcomes. If \(P_O(m)\) is its observable probability law and \(\tau(m)\) the target, its identified set at observable law \(P\) is
\[
 \mathcal I_\tau(P)=\{\tau(m):m\in\mathcal M,\ P_O(m)=P\}.
\]
Point identification means that this set is a singleton. An empty set signals incompatibility of data and assumptions. Coordinate bounds need not describe the joint set. Bounds are sharp only if every retained target is attainable or arbitrarily approachable by compatible models. Finite bounds require explicit restrictions on unobserved quantities; unrestricted confounding, hidden wealth, extrapolation or arbitrary equilibrium selection can yield uninformative sets.

\textbf{Causal interventions.} A potential outcome \(Y_i(p)\) is the outcome under a complete policy regime \(p\), including timing, financing, information and market spillovers. The target population and units are fixed before treatment. With interference, use the complete assignment vector or a justified exposure mapping; individual no-interference notation alone cannot identify an economy-wide intervention. A difference of mean potential outcomes does not require the cross-world joint distribution. Conditioning on post-treatment migration, survival, enrollment or employment changes the estimand and needs separate justification.

\textbf{Statistical statements.} An estimator, confidence set, forecast or test specifies a sampling law, model class and loss/coverage criterion. Uniform \((1-\alpha)\)-coverage means \(\inf_{m\in\mathcal M}\Pr_m\{\tau(m)\in C_n\}\geq1-\alpha\), or an explicitly asymptotic version. The whole parameter space trivially has coverage, so informativeness requires a stated width, risk or power objective. A forecasting improvement is relative to a named benchmark, information set and evaluation population.

\textbf{Equilibrium and optimization.} An equilibrium comprises feasible plans, optimality, beliefs where needed, budget constraints and all market/resource clearing. The primitives or a stated benchmark define each correspondence \(\mathcal E(p;m)\). Nonemptiness is assumed or proved, never inferred just from compact policy sets. A selected-equilibrium target uses a declared selection rule; a robust target quantifies over all permitted equilibria. Use suprema or infima unless attainment is established. An \(\varepsilon\)-optimal policy is feasible and within \(\varepsilon\) of the finite optimum value. Report an empty feasible set separately.

\textbf{Welfare and accounting.} Normative utility, interpersonal normalization, social weights, numeraire, discounting and the use of public receipts are primitives. Behavioral utility need not coincide with the welfare criterion. Household welfare includes ownership income and rents once; adding profits again to compensating variation generally double-counts them. A monetary equivalent is defined by an explicit transfer and price convention, with monotonicity and range conditions. Average effects, mechanism contributions, transfers and welfare losses are different targets. Infinite sums require convergence and dynamic feasible plans require terminal or no-Ponzi conditions.

\textbf{Source versions.} A working-paper date, revision date and journal publication year are distinguished. A DOI for a journal version is not automatically the DOI of an earlier manuscript. Locators refer to the stated version and printed pages where specified. A metadata-only check does not verify a claim about a full-text conclusion. Every entry's source subsection narrows the provenance claim accordingly.
% END ECONOMICS STATEMENT
\end{document}
